% CMP_SC 8050 progress checkpoint (about a page and a half). Standalone:
% compile this file on its own. Written to be read without the proposal deck,
% the problem statement or the report. Sources are numbered by hand in order
% of first citation; * marks those added since the proposal (Oct 2026). Not
% compiled on the build host (no TeX): check the length on Overleaf.
% Evidence as of 20 Nov 2026: commits 6a2205c5, 28cd249a, 199f03ad, 32c3a32f;
% the 58 tests in tests/unit/test_cs8050_*.py pass (19 plan checker, 39
% reductions). Parts 2 and 3 of the proof outline are Lemmas 1 and 2 of
% DeveloperDocs/cs8050/report/t1_hardness.tex; the problem is
% DeveloperDocs/cs8050/problem_statement.md; the experiment grid is steps 5.1
% to 5.9 of DeveloperDocs/FeRRy_CS8050_Project_Plan.md.
% Status gives no course due date, by choice (20 Nov 2026).
\documentclass[10pt]{article}

\usepackage[margin=0.7in]{geometry}
\usepackage{amsmath,amssymb}
\usepackage{booktabs}
\usepackage{tabularx}
\usepackage{enumitem}

\setlist{nosep,leftmargin=*}
\setlength{\parindent}{0pt}
\setlength{\parskip}{1pt}
\pagestyle{empty}

\newcommand{\PLAN}{\textsc{Plan}}
\newcommand{\head}[1]{\par\vspace{3pt}\noindent\textbf{#1}\enspace}

\begin{document}

{\large\textbf{CMP\_SC 8050 Progress Checkpoint: how hard is FeRRy's trip
planning, and how close does its planner get to the best plan?}}\\
K.~Kostage (solo project) \hfill 20 November 2026

\head{The project.} FeRRy [1] is a federated-learning system for devices with
no network: each device trains a shared model on its own data, and a drone
flies from a dock, collects their model updates and carries them back. Before
each flight, its \emph{planner} chooses one of three radio settings and a
route. The wide-band setting sends fast but reaches about 60\,m; the
narrow-band one sends slowly but reaches about 232\,m (a medium one lies
between), so the setting decides where the drone must stop. The planner tries
every plan only up to 6 devices. Beyond that it uses two heuristics with no
guarantee: a \emph{grouping rule} that places stops (one stop for the most
urgent device and every device in reach of it, repeated) and a \emph{route
search} that starts from a short loop through the stops and improves it by
small changes.
\textbf{Q1:} how hard is the planning problem? (\textbf{T1}, a hardness
proof.)
\textbf{Q2:} how far is the planner from the best plan at 7 to 20 devices?
(\textbf{T2}, an exact solver for small cases, and \textbf{T3}, worst-case
bounds on the two heuristics.)
\textbf{Q3:} for which numbers of devices, field sizes and data sizes does the
long-reach setting give the better plan? (\textbf{T4}, a formula for that
tipping point.)

\head{The problem, \PLAN.} A simplified form of what the planner solves,
checked against its code. \emph{Input:} $n$ devices, each with a position, an
importance (the value of serving it), a deadline and a data size; radio
settings, each with a reach and a sending rate that steps down as a device
gets farther from the drone; a list of \emph{sites} where the drone may stop
(for FeRRy, the stops its grouping rule proposes); the dock, the drone's
speed, a time limit $L$ and a target $W$. \emph{A plan} chooses one setting
and an ordered list of stops; each stop is a site and a chosen, non-empty set
of devices within reach of it, and no device is served twice. The drone leaves
the dock at time $0$. Devices at a stop send one after another, each for its
data size divided by its rate; flying takes distance divided by speed. \emph{A
plan is valid} if every stop finishes by the deadline of each device it serves
and the drone is back at the dock by time $L$. \PLAN{} asks whether a valid
plan serves total importance at least $W$.

\head{Status.} \emph{Completed since the proposal:} phases 1 and 2 in the
table. Q1 is answered: T1's proof is written in full in the report draft and
outlined below. \emph{Remaining:} phases 3 to 5, which answer Q2 and Q3; none
is started. \emph{Updated plan:} work pauses until a paper deadline on 30
November, then about three weeks of work follow, in the table's
order.\par\vspace{2pt}
{\small\noindent
\begin{tabularx}{\linewidth}{@{}lXll@{}}
  \toprule
  Phase & Work & Done when & Status \\
  \midrule
  1. Problem & \PLAN{} matched to the planner's code; a plan checker
    & Right on 5 hand-built plans & \textbf{Done} \\
  2. Hardness (T1) & The two reductions below; easy cases; reductions coded
    & Proved in both directions & \textbf{Done} \\
  \midrule
  3. Exact solver (T2) & Test layouts; exact solver; planner run code; upper
    bound & Equals brute-force search & Not started \\
  4. Bounds (T3, T4) & Grouping bound; route-search bad cases; tipping point
    & Proved, or counterexample & Not started \\
  5. Experiments & Gap study (below); timing; one-heuristic-off runs
    & One command reruns all & Not started \\
  \bottomrule
\end{tabularx}}\par\vspace{2pt}
\emph{Scope changes since the proposal} (the reasons are under Problems
below): the routing reduction starts from Hamiltonian cycle on grid graphs,
not the traveling salesman problem; NP-completeness is claimed only when
flying times are rounded; a deadline now binds the end of the whole stop, not
the drone's arrival; and a setting's sending rate steps down with distance, as
in the planner's code, where the proposal had one rate per setting.
\emph{Fallbacks:} if T2 cannot solve cases up to 20 devices, the sizes it
misses use the upper bound; if T3 finds no provable bound on the route search,
the report gives the largest gap found.

\head{Evidence of progress.}
\begin{itemize}
  \item \textbf{Proof outline (T1).} \emph{Claim:} \PLAN{} is strongly
    NP-hard, and strongly NP-complete when each leg's flying time is rounded
    up to a multiple of a fixed step $\delta$. \emph{Strongly} means hard
    even when every number in the input is at most polynomial in $n$, so
    unless P${}={}$NP there is no pseudo-polynomial algorithm (one polynomial
    in the numbers' values, like knapsack's dynamic program). Parts 2 and 3
    are reductions: each converts any instance of a known NP-complete problem
    into an instance of \PLAN{} with the same answer.
    \begin{enumerate}
      \item \emph{Plans are quick to check, so \PLAN{} with rounded times is
        in NP.} A plan lists at most $n$ stops and their devices. Reach is
        tested by comparing squared distances, and rounded flying times make
        every arrival and finishing time a rational number, compared exactly. Exact
        times would need comparing sums of square roots, with no known
        polynomial method [2], so NP-completeness is for rounded times only.
      \item \emph{Choosing the route alone is NP-hard: a reduction from
        Hamiltonian cycle on grid graphs.} A grid graph is a finite set of points with
        whole-number coordinates, joined when exactly $1$ apart; whether a
        cycle visits every point once is NP-complete [3]. Put a
        device with importance $1$, no data and no deadline on each of the
        $n$ points; make the points the sites and one of them the dock; use
        one setting of reach $\tfrac12$, so a device is served only from its
        own point; let a unit of distance take one time unit (equal to
        $\delta$ when times are rounded); set $L=n$ and $W=n$. A leg between
        joined points takes time $1$; points not joined are at least
        $\sqrt2$ apart, so a leg between them takes longer ($2$ once rounded
        up). A plan worth $n$ visits all $n$ points and returns, so it has at
        least $n$ legs; to finish by time $n$ every leg must take time $1$,
        so the route is a Hamiltonian cycle. Conversely, a Hamiltonian cycle
        flown from the dock is a valid plan worth $n$. All numbers are at
        most $n$, so this hardness is strong.
      \item \emph{Choosing the devices alone is NP-hard: a reduction from
        knapsack}
        (do some items, each with a size and a value, fit a capacity and
        reach a target value? NP-complete [4]). Make the dock the only site,
        with every device in reach. Each item becomes a device that sends
        for the item's size and has its value as importance; $L$ is the
        capacity and $W$ the target value. The drone never flies, so valid
        plans are exactly the sets of items that fit. This hardness is only
        weak (knapsack has a pseudo-polynomial algorithm), and it needs many
        different sending times. FeRRy has few: every device sends the same
        amount of data, and a setting has at most 15 rates, so one stop is
        solved in polynomial time by trying every count of devices per rate.
        Part 2, in contrast, still works with equal data and FeRRy's three
        settings (space the grid wider than the longest reach), so the route
        is what makes FeRRy's own planning hard.
    \end{enumerate}
    \emph{Easy special cases.} With a single site, or at most $k$ stops for a
    fixed $k$, dynamic programming solves \PLAN{} in pseudo-polynomial time
    (single site: the scheduling recursion of [5]); if sending times are also
    equal, a single site is solved greedily, as a matroid [6].
  \item \textbf{Implementation and results.} In Python: \PLAN{} with a plan
    checker that validates a plan and adds up its importance (19 tests), and
    code that builds both reductions' instances (39 tests). All 58 pass: on
    an 8-point grid graph, for example, a Hamiltonian cycle finishes exactly
    at time $L$ and a route with one diagonal leg is rejected. The planner is
    not measured yet; that needs T2.
  \item \textbf{Experiment design} (Q2, Q3). Gap $=1-{}$importance served by
    the planner\,/\,importance served by the best plan. Sizes: 6, 7, 8, 10,
    12, 14, 16, 18 and 20 devices against T2 (6 is a baseline, where the
    planner still tries everything), and 24, 48 and 96 against a
    linear-programming upper bound on the best plan, which can only overstate
    the gap. The best plan is found over two site lists: the planner's own
    stops, which shows what the route search loses, and a fuller list adding
    every device position and every crossing of two devices' reach circles,
    which also shows what the grouping rule loses. Layout families: spread
    out, clustered, or constructed to defeat the heuristics; fields of 200,
    300 and 500\,m; a loose or a tight time limit. In all, $12$ sizes
    $\times\,3\times3\times2\times30$ random layouts $=6{,}480$ cases, with
    95\% confidence intervals, reruns with one heuristic off, and, for Q3,
    the setting the planner picks against the one T4 predicts.
  \item \textbf{Literature.} With no data or deadlines \PLAN{} is
    orienteering (most value on a round trip within a time limit) [7],
    NP-hard with constant-factor approximations [8]; part 2 shows the
    hardness survives straight-line travel times. If the drone could stop
    anywhere, \PLAN{} would be a TSP with neighbourhoods [9], which no exact
    method solves at scale; hence the site list. T2 follows the Held--Karp
    dynamic program [10]; T4 uses a formula for the length of a shortest
    tour through random points [11].
\end{itemize}

\head{Problems and failed approaches.} (1) The proposal planned to reduce from
the Euclidean traveling salesman problem [12]; rounding lengthens different
tours by different amounts, which breaks it, whereas the grid's gap between
$1$ and $\sqrt2$ survives. (2) Membership in NP could not be shown
for exact times (part 1). (3) The knapsack reduction does not cover FeRRy's
own instances (part 3). (4) The proposal required the drone to arrive at a device
before its deadline; the planner's code requires the whole stop to finish by
then, and \PLAN{} now follows the code. (5) Open: the proofs are complete on
paper, but a machine check that no-instances convert to no-instances needs
T2.

\head{Sources} ($^*$ added since the proposal).
{\footnotesize
[1] Kostage et al., FeRRy, NOMS 2027 draft.
[2]$^*$ Garey, Graham, Johnson, STOC 1976.
[3]$^*$ Itai, Papadimitriou, Szwarcfiter, \emph{SIAM J. Comput.} 11(4), 1982.
[4] Garey, Johnson, \emph{Computers and Intractability}, 1979.
[5]$^*$ Lawler, Moore, \emph{Management Sci.} 16(1), 1969.
[6]$^*$ Cormen et al., \emph{Introduction to Algorithms}, 3rd ed., \S16.5.
[7] Golden, Levy, Vohra, \emph{Naval Res. Logistics} 34(3), 1987.
[8] Blum et al., \emph{SIAM J. Comput.} 37(2), 2007.
[9] Dumitrescu, Mitchell, \emph{J. Algorithms} 48(1), 2003.
[10] Held, Karp, \emph{J. SIAM} 10(1), 1962.
[11] Beardwood, Halton, Hammersley, \emph{Proc. Camb. Phil. Soc.} 55(4), 1959.
[12] Papadimitriou, \emph{Theor. Comp. Sci.} 4(3), 1977.\par}

\end{document}
